\documentclass[11pt]{article}

\usepackage[margin=1in]{geometry}
\usepackage{amsmath,amssymb,amsthm}
\usepackage{algorithm}
\usepackage{algpseudocode}
\usepackage{enumitem}
\usepackage{hyperref}
\usepackage{xcolor}
\usepackage{tikz}
\usepackage{booktabs}

\hypersetup{colorlinks=true, linkcolor=blue, urlcolor=blue, citecolor=blue}

\newtheorem{theorem}{Theorem}[section]
\newtheorem{lemma}[theorem]{Lemma}
\newtheorem{proposition}[theorem]{Proposition}
\newtheorem{corollary}[theorem]{Corollary}
\newtheorem{definition}[theorem]{Definition}
\newtheorem{example}[theorem]{Example}

\title{TOAE209/TOAH209 -- Design and Analysis of Algorithms\\[4pt]
\large Week 15: Randomized Algorithms}
\author{Foreign Trade University}
\date{}

\begin{document}
\maketitle
\tableofcontents

\section{Introduction: Flipping Coins to Compute}

Every algorithm we have studied so far has been \textbf{deterministic}: given the
same input, it always performs the same steps and produces the same answer. This
week (following Kleinberg \& Tardos, Ch.\ 13) we add a new and surprisingly
powerful ingredient -- \textbf{randomness}. A \emph{randomized algorithm} is
allowed to flip coins (call a random number generator) during its execution, and
its behaviour on a fixed input may differ from run to run.

Why would we want this? Three recurring reasons:

\begin{itemize}
    \item \textbf{Simplicity and speed.} Randomized QuickSort and randomized
    Quickselect are shorter and, on average, as fast as or faster than their best
    deterministic counterparts, while sidestepping worst-case adversarial inputs.
    \item \textbf{Breaking symmetry.} When many identical processes must make a
    decentralized decision (contention resolution), randomness lets each break the
    tie independently -- something deterministic identical processes cannot do.
    \item \textbf{Solving problems we don't know how to solve deterministically
    (efficiently).} Karger's global min-cut algorithm and the Miller--Rabin
    primality test are strikingly simple randomized algorithms for problems whose
    deterministic solutions are far more involved.
\end{itemize}

The price is that we must reason about \emph{probability}: an analysis no longer
ends with ``the answer is correct in time $T(n)$,'' but with statements like ``the
\emph{expected} running time is $O(n \log n)$'' or ``the answer is correct with
probability at least $1 - \delta$.'' The mathematics is elementary -- expectation,
linearity of expectation, indicator variables, and the union bound -- but using it
fluently is the real skill this week develops.

\subsection{Two species: Las Vegas and Monte Carlo}

\begin{definition}[Las Vegas algorithm]
A \emph{Las Vegas} algorithm is \textbf{always correct}, but its running time is a
random variable. We analyze its \emph{expected} running time. (Randomized
QuickSort and Quickselect are Las Vegas algorithms.)
\end{definition}

\begin{definition}[Monte Carlo algorithm]
A \emph{Monte Carlo} algorithm runs in a fixed (or bounded) time, but may be
\textbf{wrong with some bounded probability}. We analyze its \emph{success
probability}, and usually \emph{amplify} it by repetition. (Karger's min-cut and
Miller--Rabin are Monte Carlo algorithms.)
\end{definition}

A useful slogan: \emph{Las Vegas gambles with time; Monte Carlo gambles with
correctness.} Problem~9 asks you to implement one of each for the same task and
contrast their guarantees directly. Note that a Monte Carlo algorithm with a
\emph{checkable} answer can often be turned into a Las Vegas one (repeat until the
check passes).

\section{Probability Review}

We work in a finite (or countable) probability space. The only tools we need are
the following.

\begin{definition}[Expectation]
For a discrete random variable $X$ taking values $x$ with probabilities
$\Pr[X = x]$, the \emph{expectation} (mean) is
\[
    \mathbb{E}[X] = \sum_{x} x \cdot \Pr[X = x].
\]
\end{definition}

\begin{definition}[Indicator random variable]
For an event $A$, the \emph{indicator} $\mathbf{1}_A$ (also written $X_A$) is $1$
if $A$ occurs and $0$ otherwise. Its expectation is exactly the probability of the
event:
\[
    \mathbb{E}[\mathbf{1}_A] = 1 \cdot \Pr[A] + 0 \cdot \Pr[\bar A] = \Pr[A].
\]
\end{definition}

The single most important tool of the week is the following, which holds
\emph{whether or not} the variables are independent.

\begin{theorem}[Linearity of expectation]
\label{thm:linearity}
For any random variables $X_1, \ldots, X_n$ (defined on the same probability
space) and constants $a_1, \ldots, a_n$,
\[
    \mathbb{E}\!\left[\sum_{i=1}^{n} a_i X_i\right] = \sum_{i=1}^{n} a_i\, \mathbb{E}[X_i].
\]
\end{theorem}

The power of linearity is that it lets us compute the expectation of a complicated
quantity (e.g.\ ``number of comparisons made by QuickSort'') by writing it as a
\emph{sum of indicators}, computing the easy probability of each indicator, and
adding up -- never needing to know how the indicators interact. This pattern --
\textbf{decompose into indicators, then sum} -- recurs in nearly every proof below.

\begin{definition}[Independence; conditional probability]
Events $A$ and $B$ are \emph{independent} if $\Pr[A \cap B] = \Pr[A]\Pr[B]$. The
\emph{conditional probability} of $A$ given $B$ (with $\Pr[B] > 0$) is $\Pr[A \mid
B] = \Pr[A \cap B] / \Pr[B]$. We will repeatedly bound a product of conditional
probabilities, as in Karger's analysis.
\end{definition}

\begin{theorem}[Union bound]
\label{thm:union}
For any events $A_1, \ldots, A_n$ (no independence required),
\[
    \Pr\!\left[\bigcup_{i=1}^{n} A_i\right] \le \sum_{i=1}^{n} \Pr[A_i].
\]
\end{theorem}

The union bound is the workhorse for proving that ``nothing bad ever happens'': if
each of $n$ bad events is individually unlikely, their union is still unlikely.
We use it to bound the probability that \emph{amplified} randomized algorithms
ever fail.

\section{Contention Resolution: Breaking Symmetry}

Our first case study is the cleanest illustration of why randomness helps at all.

\begin{definition}[Contention resolution]
There are $n$ identical processes $P_1, \ldots, P_n$ all wanting to access a single
shared resource (e.g.\ a broadcast channel). In each discrete \emph{round}, every
process may attempt to access the resource. An attempt by $P_i$ \emph{succeeds} in
a round only if $P_i$ is the \emph{only} process attempting that round; if two or
more attempt simultaneously, all of them fail (a ``collision''). The processes
cannot communicate or coordinate. How can they all eventually succeed?
\end{definition}

If the processes are deterministic and identical they are doomed: they will all do
exactly the same thing every round and collide forever. The fix is randomness:
each process attempts \emph{independently} with some probability $p$ per round.

\begin{algorithm}[h]
\caption{Randomized Contention Resolution (one process $P_i$)}
\begin{algorithmic}[1]
\While{$P_i$ has not yet succeeded}
    \State With probability $p$, attempt to access the resource this round
    \State (otherwise, stay silent this round)
\EndWhile
\end{algorithmic}
\end{algorithm}

Let $A_{i,t}$ be the event ``$P_i$ attempts and succeeds in round $t$.'' Process
$P_i$ succeeds in round $t$ exactly when it attempts (probability $p$) \emph{and}
no other process attempts (probability $(1-p)^{\,n-1}$):
\[
    \Pr[A_{i,t}] = p\,(1-p)^{\,n-1}.
\]
This is maximized at $p = 1/n$, giving
\[
    \Pr[A_{i,t}] = \frac1n\Bigl(1 - \frac1n\Bigr)^{n-1} \;\ge\; \frac{1}{e\,n},
\]
using the standard fact $(1 - 1/n)^{n-1} \ge 1/e$. So with $p = 1/n$ each process
succeeds in any given round with probability $\Theta(1/n)$.

\begin{theorem}[All processes succeed in $O(n \log n)$ rounds]
\label{thm:contention}
With $p = 1/n$, after $t = \lceil 2 e n \ln n\rceil$ rounds, all $n$ processes have
succeeded at least once with probability at least $1 - 1/n$.
\end{theorem}

\begin{proof}
Fix a process $P_i$. By the above, in each round $P_i$ fails to succeed with
probability at most $1 - \frac{1}{en}$. Rounds are independent, so the probability
$P_i$ has \emph{not} succeeded in any of the first $t$ rounds is at most
\[
    \Bigl(1 - \tfrac{1}{en}\Bigr)^{t} \le e^{-t/(en)},
\]
using $1 - x \le e^{-x}$. Setting $t = \lceil 2en\ln n\rceil$ gives this at most
$e^{-2\ln n} = 1/n^2$. Now apply the \textbf{union bound} (Theorem~\ref{thm:union})
over the $n$ ``failure'' events: the probability that \emph{some} process has not
yet succeeded is at most $n \cdot 1/n^2 = 1/n$. Hence with probability $\ge 1 -
1/n$ all $n$ processes have succeeded within $O(n \log n)$ rounds.
\end{proof}

Problem~5 asks you to simulate this and check empirically that the number of rounds
until all succeed grows like $\Theta(n \log n)$.

\section{Randomized QuickSort}

\subsection{The algorithm}

QuickSort partitions an array around a \emph{pivot} and recurses. Its worst case
is $\Theta(n^2)$ -- triggered, for a fixed pivot rule, by adversarial inputs (e.g.\
an already-sorted array with ``first element'' as pivot). The randomized version
chooses the pivot \emph{uniformly at random}, which makes the expected running time
$O(n \log n)$ on \emph{every} input, with no bad inputs at all.

\begin{algorithm}[h]
\caption{Randomized QuickSort}
\begin{algorithmic}[1]
\Function{RQuickSort}{$A$}
    \If{$|A| \le 1$} \Return $A$ \EndIf
    \State Choose a pivot index uniformly at random; let $x$ be that element
    \State Partition $A \setminus \{x\}$ into $L = \{y < x\}$ and $R = \{y > x\}$
    \State \Return $\textsc{RQuickSort}(L) \mathbin{+} [x] \mathbin{+} \textsc{RQuickSort}(R)$
\EndFunction
\end{algorithmic}
\end{algorithm}

\subsection{Expected number of comparisons}

The running time is dominated by element comparisons, so we count those.

\begin{theorem}
\label{thm:quicksort}
The expected number of comparisons made by randomized QuickSort on an array of $n$
distinct elements is at most $2n \ln n = O(n \log n)$.
\end{theorem}

\begin{proof}
Let the sorted order of the elements be $z_1 < z_2 < \cdots < z_n$. The whole
analysis rests on one clean observation: \emph{any two elements $z_i$ and $z_j$ are
compared at most once over the entire execution} (a comparison happens only against
a pivot, and once an element is a pivot it is removed from all recursive calls).

For $i < j$, let $X_{ij}$ be the indicator of the event ``$z_i$ and $z_j$ are ever
compared.'' The total number of comparisons is $X = \sum_{i<j} X_{ij}$, so by
\textbf{linearity of expectation} (Theorem~\ref{thm:linearity}),
\[
    \mathbb{E}[X] = \sum_{i<j} \mathbb{E}[X_{ij}] = \sum_{i<j} \Pr[z_i \text{ and } z_j \text{ are compared}].
\]

\textbf{Key probability.} Consider the set $Z_{ij} = \{z_i, z_{i+1}, \ldots, z_j\}$
of the $j - i + 1$ elements between $z_i$ and $z_j$ inclusive. Track the first
element of $Z_{ij}$ ever chosen as a pivot (every element of $Z_{ij}$ is eventually
a pivot or separated from the others only by some element of $Z_{ij}$ being chosen
first). Two cases:
\begin{itemize}
    \item If the first pivot chosen from $Z_{ij}$ is $z_i$ or $z_j$, then $z_i$ and
    $z_j$ are compared (the pivot is compared to everyone in its subarray, which
    still contains the other).
    \item If the first pivot chosen from $Z_{ij}$ is some $z_k$ with $i < k < j$,
    then $z_i$ and $z_j$ are split into different subarrays and \emph{never}
    compared.
\end{itemize}
Since the pivot is chosen uniformly at random, the first pivot from $Z_{ij}$ is
equally likely to be any of its $j - i + 1$ elements. Exactly $2$ of those
($z_i$ and $z_j$) cause a comparison, so
\[
    \Pr[z_i \text{ and } z_j \text{ compared}] = \frac{2}{j - i + 1}.
\]

\textbf{Summing.} Therefore
\[
    \mathbb{E}[X] = \sum_{i=1}^{n-1}\sum_{j=i+1}^{n} \frac{2}{j-i+1}
    \le \sum_{i=1}^{n} \sum_{k=2}^{n} \frac{2}{k}
    = 2n \sum_{k=2}^{n} \frac1k
    \le 2n \ln n,
\]
where the last step uses the harmonic bound $\sum_{k=1}^{n} 1/k \le 1 + \ln n$.
Hence $\mathbb{E}[X] = O(n \log n)$.
\end{proof}

Problem~2 asks you to implement randomized QuickSort and verify both correctness
(against \texttt{sorted()}) and that the comparison count tracks $O(n \log n)$.

\section{Randomized Quickselect}

The same idea, restricted to one side of the recursion, finds the $k$-th smallest
element (the $k$-th \emph{order statistic}) in expected \emph{linear} time.

\begin{algorithm}[h]
\caption{Randomized Quickselect (find the $k$-th smallest)}
\begin{algorithmic}[1]
\Function{RSelect}{$A, k$}
    \If{$|A| = 1$} \Return $A[0]$ \EndIf
    \State Choose a pivot $x \in A$ uniformly at random
    \State Partition into $L = \{y < x\}$ and $R = \{y > x\}$
    \If{$k \le |L|$} \Return $\textsc{RSelect}(L, k)$
    \ElsIf{$k = |L| + 1$} \Return $x$
    \Else{} \Return $\textsc{RSelect}(R,\, k - |L| - 1)$
    \EndIf
\EndFunction
\end{algorithmic}
\end{algorithm}

\begin{theorem}
\label{thm:quickselect}
Randomized Quickselect runs in expected $O(n)$ time on an array of size $n$.
\end{theorem}

\begin{proof}[Proof sketch]
Call a pivot \emph{good} if it lands in the middle half (between the $25$th and
$75$th percentiles); a uniformly random pivot is good with probability $\ge 1/2$. A
good pivot shrinks the surviving subarray to at most $\tfrac34$ of its size.
Group the recursion into ``phases,'' where a phase ends each time a good pivot is
hit. The expected number of pivots per phase is at most $2$ (a geometric random
variable with success probability $\ge 1/2$), and the subarray size shrinks by a
factor $\ge \tfrac34$ each phase. So the expected total work is at most
\[
    \sum_{\text{phases } j \ge 0} 2 \cdot c\,\Bigl(\tfrac34\Bigr)^{j} n
    = 2cn \sum_{j \ge 0}\Bigl(\tfrac34\Bigr)^{j} = 2cn \cdot 4 = O(n),
\]
since the work in a subarray of size $m$ (partitioning) is $\Theta(m)$ and the
geometric series converges. Hence the expected running time is $O(n)$.
\end{proof}

Problem~3 asks you to implement Quickselect and verify its output against
\texttt{sorted()[k-1]} over many seeded instances.

\section{Karger's Randomized Global Min-Cut}

\subsection{The problem and the algorithm}

\begin{definition}[Global minimum cut]
Given a connected undirected (multi)graph $G = (V, E)$, a \emph{global cut} is a
partition of $V$ into two non-empty sets $(S, V \setminus S)$; its \emph{size} is
the number of edges with one endpoint on each side. The \emph{global min-cut} is a
cut of minimum size.
\end{definition}

Karger's algorithm is astonishingly simple: repeatedly \emph{contract} a random
edge (merging its two endpoints into one ``super-vertex,'' keeping parallel edges
but deleting self-loops) until only two super-vertices remain. The edges between
those two final super-vertices form a cut; return it.

\begin{algorithm}[h]
\caption{Karger's Contraction Algorithm}
\begin{algorithmic}[1]
\Function{Contract}{$G = (V, E)$}
    \While{$|V| > 2$}
        \State Choose an edge $e = (u, v) \in E$ uniformly at random
        \State Contract $e$: merge $u$ and $v$ into one vertex; keep parallel
        edges; delete resulting self-loops
    \EndWhile
    \State \Return the multiset of edges between the two remaining vertices (a cut)
\EndFunction
\end{algorithmic}
\end{algorithm}

\subsection{Success probability}

\begin{theorem}
\label{thm:karger}
A single run of Karger's contraction algorithm returns a fixed global min-cut with
probability at least $\dfrac{1}{\binom{n}{2}} = \dfrac{2}{n(n-1)}$.
\end{theorem}

\begin{proof}
Fix one particular global min-cut $C^*$ of size $k$. The algorithm returns $C^*$
exactly when it \emph{never} contracts any edge of $C^*$. We bound the probability
of this surviving sequence of contractions.

First, note that every vertex must have degree $\ge k$: if some vertex $v$ had
degree $< k$, the cut $(\{v\}, V\setminus\{v\})$ would be smaller than $k$,
contradicting minimality. So at any stage with $m$ super-vertices, the total number
of edges is $\ge mk/2$ (sum of degrees $\ge mk$, divided by $2$).

At the $i$-th contraction step there are $n - i + 1$ super-vertices, hence at least
$(n-i+1)k/2$ edges, of which exactly $k$ belong to $C^*$. So the probability of
contracting a $C^*$-edge at this step is at most
\[
    \frac{k}{(n-i+1)k/2} = \frac{2}{n-i+1}.
\]
The probability we \emph{avoid} all $C^*$-edges through all $n-2$ contractions is,
by the chain rule,
\[
    \Pr[C^* \text{ survives}] \ge \prod_{i=1}^{n-2}\Bigl(1 - \frac{2}{n-i+1}\Bigr)
    = \prod_{i=1}^{n-2}\frac{n-i-1}{n-i+1}.
\]
This telescopes:
\[
    \frac{n-2}{n}\cdot\frac{n-3}{n-1}\cdot\frac{n-4}{n-2}\cdots\frac{2}{4}\cdot\frac{1}{3}
    = \frac{2}{n(n-1)} = \frac{1}{\binom{n}{2}}. \qedhere
\]
\end{proof}

\subsection{Amplification}

A success probability of $1/\binom{n}{2} \approx 2/n^2$ seems tiny, but it is
\emph{polynomially} small, so a polynomial number of independent repetitions makes
failure exponentially unlikely.

\begin{corollary}[Amplification]
\label{cor:karger-amp}
Running Karger's algorithm $T = \binom{n}{2}\ln n$ times independently and keeping
the smallest cut found returns a global min-cut with probability at least $1 -
1/n$.
\end{corollary}

\begin{proof}
Each run independently fails to find the fixed min-cut $C^*$ with probability at
most $1 - 1/\binom{n}{2}$. The probability all $T$ runs fail is at most
\[
    \Bigl(1 - \tfrac{1}{\binom{n}{2}}\Bigr)^{T} \le e^{-T/\binom{n}{2}} = e^{-\ln n} = \frac1n,
\]
using $1 - x \le e^{-x}$. So with probability $\ge 1 - 1/n$ at least one run finds
$C^*$, and keeping the best cut returns it.
\end{proof}

Problem~6 implements a single contraction run; Problem~7 adds amplification and
checks against a brute-force / max-flow min-cut on small graphs; Problem~11
empirically verifies the $1/\binom{n}{2}$ floor on the single-run success
probability.

\section{Reservoir Sampling}

Often we must sample from a \emph{stream} of unknown length -- we see items one at
a time and cannot store them all. \emph{Reservoir sampling} maintains a uniform
random sample of $k$ items using only $O(k)$ space and a single pass.

\begin{algorithm}[h]
\caption{Reservoir Sampling (sample $k$ of an $n$-item stream)}
\begin{algorithmic}[1]
\State $R \gets$ the first $k$ items of the stream
\For{$t = k+1, k+2, \ldots$ (each new item $x_t$)}
    \State With probability $k/t$, replace a uniformly random element of $R$ with $x_t$
\EndFor
\State \Return $R$
\end{algorithmic}
\end{algorithm}

\begin{theorem}
\label{thm:reservoir}
After processing a stream of $n \ge k$ items, every item is in the reservoir $R$
with probability exactly $k/n$; equivalently, $R$ is a uniform random $k$-subset.
\end{theorem}

\begin{proof}
We show item $x_t$ ends in $R$ with probability $k/n$, for every $t$. For $t \le
k$, $x_t$ starts in $R$ (probability $1$); for $t > k$, $x_t$ enters $R$ with
probability $k/t$. In either case, after step $t$ the item is in $R$ with
probability $k/t$ (for $t \le k$ this is $1 \ge$, handled directly; we focus on the
inductive survival). For each later step $s > t$, $x_t$ \emph{survives} unless step
$s$ both replaces (probability $k/s$) and selects $x_t$'s slot (probability
$1/k$), i.e.\ it is removed with probability $1/s$, surviving with probability
$1 - 1/s = (s-1)/s$. Hence
\[
    \Pr[x_t \in R \text{ at the end}] = \frac{k}{t}\prod_{s=t+1}^{n}\frac{s-1}{s}
    = \frac{k}{t}\cdot\frac{t}{n} = \frac{k}{n},
\]
the product telescoping. (For $t \le k$ the leading factor is $1$ and the same
telescoping gives $k/n$.) Since every item is included with the same probability
$k/n$, $R$ is a uniform random $k$-subset.
\end{proof}

Problem~4 asks you to implement reservoir sampling and verify the $k/n$ inclusion
frequency empirically over many seeded trials.

\section{Hashing and Number-Theoretic Randomized Algorithms}

\subsection{Universal hashing}

A \emph{hash function} maps keys to a small range $\{0, \ldots, m-1\}$ of buckets.
A fixed hash function can always be defeated by an adversary who piles all keys
into one bucket. \emph{Universal hashing} defeats the adversary by choosing the
hash function \emph{at random} from a carefully designed family.

\begin{definition}[Universal family]
A family $\mathcal H$ of hash functions $h : U \to \{0,\ldots,m-1\}$ is
\emph{universal} if for every pair of distinct keys $x \ne y$,
\[
    \Pr_{h \in \mathcal H}[h(x) = h(y)] \le \frac1m.
\]
\end{definition}

A standard universal family: pick a prime $p > |U|$, choose $a \in \{1,\ldots,p-1\}$
and $b \in \{0,\ldots,p-1\}$ uniformly at random, and set
\[
    h_{a,b}(x) = \bigl((a x + b) \bmod p\bigr) \bmod m.
\]
Using indicator variables and linearity of expectation, the \emph{expected} number
of colliding pairs when hashing $n$ keys into $m$ buckets is at most
\[
    \binom{n}{2}\cdot\frac1m \approx \frac{n^2}{2m},
\]
since each of the $\binom n2$ pairs collides with probability $\le 1/m$.
Problem~12 asks you to implement this $(ax+b \bmod p)\bmod m$ family and verify the
expected collision count empirically.

\subsection{Monte Carlo estimation of $\pi$}

A toy but vivid Monte Carlo algorithm: throw $N$ random darts uniformly into the
unit square $[0,1]^2$ and count the fraction landing inside the quarter-circle of
radius $1$. That fraction approaches the area $\pi/4$, so $4 \times$ (fraction)
estimates $\pi$. The error shrinks like $1/\sqrt N$ (by the variance of a sum of
$N$ i.i.d.\ indicators). Problem~8 implements this and checks the estimate is
within a loose tolerance for large $N$.

\subsection{Miller--Rabin primality testing}

Deciding whether a large integer is prime by trial division costs $\Theta(\sqrt n)$
-- hopeless for cryptographic-size numbers. The \textbf{Miller--Rabin} test is a
Monte Carlo algorithm that, for a candidate $n$, picks random ``witnesses'' $a$ and
checks an algebraic condition (derived from Fermat's little theorem, refined to
catch Carmichael numbers). If $n$ is prime, the test \emph{always} reports
``probably prime.'' If $n$ is composite, each random witness exposes it with
probability $\ge 3/4$, so $r$ independent witnesses reduce the error to $\le
(1/4)^r$ -- amplification again. Problem~10 implements Miller--Rabin and checks it
against deterministic trial division over a range.

\begin{algorithm}[h]
\caption{Miller--Rabin (one witness $a$ for odd $n > 2$)}
\begin{algorithmic}[1]
\State Write $n - 1 = 2^{s} d$ with $d$ odd
\State $x \gets a^{d} \bmod n$
\If{$x = 1$ or $x = n-1$} \Return ``probably prime'' \EndIf
\For{$r = 1$ to $s-1$}
    \State $x \gets x^{2} \bmod n$
    \If{$x = n-1$} \Return ``probably prime'' \EndIf
\EndFor
\State \Return ``composite''
\end{algorithmic}
\end{algorithm}

\section{Summary and Course Wrap-Up}

\subsection{This week}

Randomization gives us a third design paradigm, alongside the deterministic
techniques (greedy, divide-and-conquer, dynamic programming, network flow) of
earlier weeks. The recurring analytical patterns were:

\begin{itemize}
    \item \textbf{Indicator variables + linearity of expectation}: count a
    complicated quantity by decomposing it into a sum of simple $0/1$ events and
    summing their probabilities -- used for QuickSort's comparisons
    (Theorem~\ref{thm:quicksort}), reservoir sampling, and universal-hashing
    collisions.
    \item \textbf{The union bound}: prove ``nothing bad ever happens'' by summing
    the small probabilities of individual bad events -- used for contention
    resolution (Theorem~\ref{thm:contention}) and Karger amplification
    (Corollary~\ref{cor:karger-amp}).
    \item \textbf{Amplification}: turn a small-but-polynomial success probability
    into an overwhelming one by independent repetition -- Karger's min-cut and
    Miller--Rabin.
    \item \textbf{Las Vegas vs.\ Monte Carlo}: gamble with time (always correct) or
    with correctness (always fast), and convert between them when the answer is
    checkable.
\end{itemize}

\subsection{Looking back over the course}

This is the final week of TOAE209/TOAH209. Over fifteen weeks we built a toolkit
for \emph{designing} and \emph{proving correct} algorithms:

\begin{itemize}
    \item \textbf{Foundations} (Weeks 1--3): stable matching, asymptotic analysis,
    graphs and union-find.
    \item \textbf{Greedy} (Weeks 4--5): exchange arguments and ``greedy stays
    ahead,'' applied to scheduling, shortest paths, and minimum spanning trees.
    \item \textbf{Divide and conquer} (Weeks 6--7): recurrences, the Master
    Theorem, mergesort, fast multiplication, and the FFT.
    \item \textbf{Dynamic programming} (Weeks 8--9): optimal substructure,
    memoization, and a gallery of DP formulations.
    \item \textbf{Network flow} (Weeks 10--11): max-flow/min-cut and its many
    reductions (bipartite matching, project selection, image segmentation).
    \item \textbf{Complexity and coping with hardness} (Weeks 12--14):
    NP-completeness, PSPACE, parameterized/structural tractability, approximation
    algorithms, and local search.
    \item \textbf{Randomization} (Week 15): a new paradigm that trades certainty
    for simplicity, speed, and the ability to break symmetry.
\end{itemize}

The unifying lesson is that an algorithm is not finished when it runs -- it is
finished when you can \emph{prove} what it does and how fast. Whether the proof is
an exchange argument, an induction on subproblems, a flow-conservation argument, an
NP-completeness reduction, or an expectation computed with indicator variables, the
discipline of pairing every algorithm with its proof is what distinguishes
algorithm design from mere programming. Carry that discipline forward.

\end{document}
